\section{Điều bài báo chưa giải quyết}
\label{sec:ch11-dieu-chua-giai-quyet}

Mục trước dừng ở một chỗ hở mà chính các tác giả ghi lại, và phần hạn chế của
bài còn ba chỗ nữa. Chỗ thứ nhất là giả thiết: toàn bộ lý thuyết dựng trên phép
che đặc trưng bằng baseline, mà chương~\ref{ch:attribution-theo-gradient} đã
cho thấy baseline là một lựa chọn tự do chứ không phải một hằng số của bài
toán. Chỗ thứ hai là chi phí: giá trị Shapley tính chính xác được với trò chơi dưới
20 người chơi, còn với trò chơi lớn hơn, chẳng hạn lời nhắc của mô hình ngôn
ngữ, thì phải dùng thuật toán xấp xỉ.
Chỗ thứ ba là phía người đọc: các tác giả viết rằng đo hiệu ứng bậc cao thì
đồng thời làm tăng chi phí tính toán và gánh nặng nhận thức cho người
dùng~\autocite{p23metagame}.

Ba chỗ ấy cộng với chỗ hở của mục trước là toàn bộ những gì bài tự nhận. Sách
thêm vào hai điều, và cả hai đã được nêu ở trên nên chỉ nhắc lại tên:
\tn{tính hiệu quả phân tầng}{hierarchical efficiency} ràng buộc quan hệ giữa
hai tầng mà không xác nhận tầng dưới, và
ba cái nền của phần thí nghiệm đều nằm ở phía tính hợp lý.

Với ngần ấy, câu hỏi mà chương~\ref{ch:lime-nao-dang-tin} giao lại đã có đủ dữ
kiện để trả lời. Câu hỏi ấy là liệu vòng của hình~\ref{fig:ch10-vong-bien-the}
có lặp lại ở những nhánh ngoài LIME hay không, và bài báo của chương này là một
điểm dữ liệu ở một nhánh xa: không phải \tn{mô hình thay thế cục bộ}{local surrogate model}, không phải
attribution đặc trưng bậc nhất, mà một tầng dựng bên trên cả hai.

Vòng ấy lặp lại, và lặp lại với từng bước một. Có một chỗ yếu đã biết của các
phương pháp sẵn có, ở đây là chỗ rò của
mục~\ref{sec:ch11-ro-ri-tuong-tac}. Có một xây dựng mới sửa đúng chỗ yếu ấy, ở
đây là meta-attribution, và nó sửa thật: bảng~\ref{tab:ch11-tach-hieu-ung} là
một chứng cứ phân tích chứ không phải một lời tuyên bố. Có phép so với những gì
có sẵn, theo những chỉ số mà người đề xuất chọn, ở đây là Recall@K trên nhãn do
chính nhóm tác giả gán, chỉ số nhận diện tương tác, và điểm phân đoạn. Và rồi
công trình được công bố.

Bước còn thiếu vẫn là bước ở mục~\ref{sec:ch10-khong-co-cach-so-sanh}, chỉ khác
là lần này nó thiếu ở một chỗ dễ nhìn hơn nhiều. Không có bước nào trong vòng
nói được rằng một phương pháp cũ nên thôi dùng. Phép so sánh làm được điều đó
là phép so sánh mà bài nêu tên rồi để lại, tức là phép đo độ trung thực; và
việc nó bị để lại không phải chuyện thiếu công sức, vì thí nghiệm đã chạy tám
GPU H100 trong khoảng 15 ngày rồi chừng ấy nữa cho những lần thất bại. Cái
thiếu không phải tài nguyên mà là một dụng cụ đo mà cả hai bên cùng chấp nhận.

Cho nên hình dạng mà chương~\ref{ch:lime-nao-dang-tin} nghi ngờ đúng là hình
dạng chung, và nó chung không vì người trong nghề làm sai điều gì. Nó chung vì
một vòng nghiên cứu chỉ loại bỏ được một phương pháp khi có chỗ để đối chiếu,
và Phần~IV đến giờ vẫn chưa tìm ra chỗ ấy. Điều chương này thêm vào là một
điểm dữ liệu ở chỗ khó chối nhất: không phải một khảo sát nhận xét về người
khác, mà một công trình tự ghi vào phần hạn chế của mình rằng phép so sánh
quyết định vẫn còn thiếu.

Mọi dụng cụ mà Phần~IV đã xét cho tới đây đều là dụng cụ không có người tham
gia. Chương~\ref{ch:con-nguoi-vang-mat} hỏi về những dụng cụ còn lại.
